\documentclass{dicomopapers}
\usepackage[utf8]{inputenc}
\usepackage[T1]{fontenc}
\usepackage{geometry}
\usepackage{graphicx}
\usepackage{float}
\usepackage{hyperref}
\usepackage[english,indonesian]{babel}

\title{Optimasi Pemrosesan Template LaTeX Menggunakan Pendekatan Berbasis AI}
\author{Academic Support Agent}
\date{\today}

\begin{document}
\maketitle

\begin{abstract}

\end{abstract}

\chapter{PENDAHULUAN}


\chapter{TINJAUAN PUSTAKA}


% Original content (commented out):
% \documentclass[Japanese]{dicomopapers}
\usepackage[dvipdfmx]{graphicx}
\usepackage{latexsym}
\usepackage{nidanfloat} % 1ページ概要内でfigure*を使うときは必須

%======================================================================
% preamble
%======================================================================
\def\Underline{\setbox0\hbox\bgroup\let\\\endUnderline}
\def\endUnderline{\vphantom{y}\egroup\smash{\underline{\box0}}\\}
\def\|{\verb|}

%概要投稿用余白調整ここから
\setlength{\Jauthorjreceivesep}{0.0mm}
\setlength{\Jreceivejabstsep}{0.0mm}
\setlength{\Jabstsepjkeyword}{0.0mm}
\setlength{\Jkeywordetitle}{0.0mm}
%概要投稿用余白調整ここまで

% 行間調整
%\def\baselinestretch{0.98}

\makeatletter
\addtolength{\JEhonbunsep}{-\baselineskip}
\setlength{\dblfloatsep}{.3\Cvs \@plus.2\Cvs \@minus.1\Cvs}
\setlength{\dbltextfloatsep}{.3\Cvs \@plus.2\Cvs \@minus.1\Cvs}
\makeatother

% bibのnewblock
\makeatletter
\def\newblock{\hskip .11em\@plus.33em\@minus.07em}
\makeatother

%----------------------------------------------------------------------
% title, authors
%----------------------------------------------------------------------
% 和文表題
\title{タイトル\\ ですよ}

% 英文表題
\etitle{This is Title}

% 所属ラベルの定義
\affiliate{IPSJ}{情報処理学会\\
Information Processing Society of Japan
}
\affiliate{HOGE}{ほげ学会\\
Hoge Society
}

\author{情報 太郎}{Taro Joho}{IPSJ}[taro@ipsj.or.jp]
\author{情報 士郎}{Shiro Joho}{IPSJ}
\author{捕下 次郎}{Jiro Hoge}{HOGE}

%======================================================================
\begin{document}
\maketitle

\small

%----------------------------------------------------------------------
\section{はじめに}

IPSJ DICOMOの概要(A4 1ページ)のためのテンプレートです．
\figref{fig:test1}〜\figref{fig:test3}は，ページ下部に図を3つ並べる例を示している．

\begin{figure*}[b]
 \begin{minipage}[t]{.32\textwidth}
  \centering
  \fbox{\rule{0pt}{3cm}\rule{3cm}{0pt}}
  \caption{図の置き方の例}
  \label{fig:test1}
 \end{minipage}
 \hfill
 \begin{minipage}[t]{.32\textwidth}
  \centering
  \fbox{\rule{0pt}{3cm}\rule{3cm}{0pt}}
  \caption{図の置き方の例}
  \label{fig:test2}
 \end{minipage}
 \hfil
 \begin{minipage}[t]{.32\textwidth}
  \centering
  \fbox{\rule{0pt}{3cm}\rule{3cm}{0pt}}
  \caption{図の置き方の例}
  \label{fig:test3}
 \end{minipage}
\end{figure*}

%----------------------------------------------------------------------
\section{}

%----------------------------------------------------------------------
\section{おわりに}

%----------------------------------------------------------------------
% references
%----------------------------------------------------------------------
%\bibliographystyle{ipsjunsrt}
%\bibliography{bib/my}

\begin{thebibliography}{1}

\bibitem{flab/ishida18:itsc}
Ishida, S.\ et~al.,: {SAVeD}:
  Acoustic Vehicle Detector with Speed Estimation capable of Sequential Vehicle
  Detection, {\em {IEEE} ITSC} (2018).

\end{thebibliography}

\end{document}


\end{document}